\chapter*{Resumen}

\section*{\tituloPortadaVal}

La invención de los \textbf{sintetizadores} en los años 60 revolucionó el mundo de la música, antes dominada por los instrumentos analógicos tradicionales. La llegada de la \textbf{música sintética} y del ordenador como elemento central de la producción moderna han definido la forma de trabajar de los artistas modificando su sonido y procesos creativos. Desde entonces los productores musicales y diseñadores de sonido han pasado años aprendiendo a dominar los sintetizadores para definir y crear timbres musicales nunca antes vistos, sin embargo, dicho proceso a menudo requiere de mucha maestría, pues los distintos tipos de síntesis suelen ser complejos, modulares y escalables, lo que lo convierte en una tarea altamente complicada para la mayoría de los músicos. Es común que incluso los mejores diseñadores de sonido encuentren dificultad en \textbf{replicar otros sonidos sintéticos}, pues dependen directamente del hardware utilizado y de la secuencia de efectos y modulaciones sonoras aplicadas.

Es por ello que la aplicación del \textbf{aprendizaje profundo} a la síntesis abre todo un abanico de posibilidades al delegar la detección de patrones y conocimiento técnico requerido a la \textbf{inteligencia artificial}, permitiendo al artista ocuparse de la experimentación y la creación que le pertenece. Mediante el uso de \textbf{Redes Neuronales Convolucionales (CNNs)}, este proyecto busca automatizar la estimación de parámetros en la compleja \textbf{síntesis FM}. Así, la IA resuelve la barrera técnica de replicar un sonido (\textbf{sound matching}), convirtiéndose en una herramienta directa al servicio de la creatividad.

\vspace{1em} 

Todo el código desarrollado se encuentra en nuestro repositorio de GitHub: \\ 
\url{https://github.com/Angelji06/TFG-SoundMatchingFM}

\section*{Palabras clave}
   
\noindent Replicación Sonora, Síntesis FM, Aprendizaje Profundo, Redes Neuronales Convolucionales, Procesamiento de Señales de Audio, Espectrogramas, Python, PyTorch.




